\clearpage\mbox{}
\chapter{Bibliography Compilation}
\label{chapter:BibliographyCompilation}

\remarksupervisor{This chapter should be commented out in the final version.} 

\begin{enumerate}
    \item Create a section for each subtopic.
    \item Using an itemize environment, write the name of the paper or book and cite it. 
    \item For each citation, using an itemize environment, make a bullet point list of the main results in the  the paper. 
\end{enumerate}


\section{Rule based trading}

\begin{enumerate}

    \item Efficient capital markets: A review of theory and empirical work.
    \cite{fama1970efficient}

    \begin{itemize}
        \item Introduces the Efficient Market Hypothesis (EMH).
        \item Argues that financial markets quickly incorporate available information.
        \item Provides theoretical foundation for evaluating the effectiveness of trading rules.
    \end{itemize}

    \item Time series momentum.
    \cite{moskowitz2012time}

     \begin{itemize}
        \item Provides empirical evidence for time-series momentum across multiple asset classes.
        \item Shows that past returns can predict future returns over intermediate horizons.
        \item Supports the theoretical foundation of trend-following strategies.
    \end{itemize}
    
    \item A Forex trading expert system based on a new approach to the rule-base evidential reasoning. 
    \cite{dymova2016forex}.
    
    \begin{itemize}
    \item They proposed rule-base evidential reasoning (REBR) the combination of fuzzy logic and the Dempster-Shafer theory of evidence
    \item They discovered that Dempster's rule may provide wrong results when there is no conflict and that the combination is key for improved results.
    \end{itemize}
    
    \item Evolving trading rules. 
    \cite{ghandar2008evolving}.
    
    \begin{itemize}
    \item Combines evolutionary computation with fuzzy logic to evolve interpretable trading rules instead of "black box" models.
    \item Uses a variable-length sliding window that automatically adjusts training periods based on real-time performance to handle market changes.
    \end{itemize}

    \item Generating trading rules on the stock markets with genetic programming. 
    \cite{potvin2004generating} 

    \begin{itemize}
        \item Applies genetic programming to automatically generate trading rules.
        \item Demonstrates that evolutionary approaches can discover profitable rule structures.
        \item Highlights the flexibility of rule-based strategy generation.
    \end{itemize}

     \item Evolving trading rule-based policies. 
    \cite{bradley2010evolving} 

    \begin{itemize}
        \item Evolves trading rule-based policies using evolutionary computation techniques.
        \item Emphasizes adaptive rule generation in dynamic market environments.
        \item Compares evolved strategies against benchmark approaches.
    \end{itemize}
    
    \item A rule-based dynamic decision-making stock trading system based on quantum-inspired tabu search algorithm 
    \cite{chou2014rule}.
    
    \begin{itemize}
        \item They propose a hybrid "Quantum-inspired Tabu Search" (QTS) algorithm to optimize the combination of technical trading rules, allowing for a more efficient search of profitable strategies compared to traditional methods.
        \item They demonstrate that the QTS-optimized system, when combined with a sliding window mechanism to prevent overfitting, significantly outperforms the buy-and-hold benchmark and other standard evolutionary approaches.
    \end{itemize}

    \item Enhanced decision making mechanism of rule-based genetic network programming for creating stock trading signals
    \cite{mabu2013enhanced}.
    \begin{itemize}
        \item Genetic Network Programming (GNP) a graph-based evolutionary algorithm to improve rule-based trading and to generate effective trading rules.
        \item Their results show that rule-based system is better than traditional individual-based trading model, but that a trend judgment mechanism is necessary to increase the chances of profit.
    \end{itemize}

    \item The Profitability Of Technical Trading
Rules: A Combined Signal Approach
    \cite{lento2007profitability}
    \begin{itemize}
        \item Investigates whether a rule-based trading approach can outperform buy and hold strategy.
        \item The result is mixed, no generalized profitable technical trading strategy.
        \item A multi-indicator strategy can still be useful for making investment decisions.
    \end{itemize}

    \item Do Technical Trading Rules Generate
Profits? Conclusions from the Intra-Day
Foreign Exchange Market
    \cite{curcio1997technical}
    \begin{itemize}
        \item Examines the profitability of technical trading rules.
        \item The result is that none of their strategies are profitable
    \end{itemize}

\end{enumerate}

\section{Indicators}

\begin{enumerate}

    \item Enhancing Trading Strategies: A Multi-indicator Analysis for Profitable Algorithmic Trading
    \cite{sukma2025enhancing}
    \begin{itemize}
        \item Compares and benchmarks indicators.
        \item Highlights the importance of multi-indicator strategies, which significantly outperformed all individual indicators.
    \end{itemize}

    \item Algorithmic Trading using Technical Indicators
    \cite{salkar2021algorithmic}
    \begin{itemize}
        \item Explores and proposes trading strategies based on quantitative analysis of time series data.
        \item The result showed that the strategy with RSI and MACD technical indicators gave the highest return of 12\%.
    \end{itemize}

    \item Trading strategies based on trading systems: Evidence from the performance of
technical indicators
    \cite{keshavarz2022trading}
    \begin{itemize}
        \item The performance of 11 technical indicators are analyzed.
        \item The result showed that creating a strategy by combining moving average, exponential moving average and relative strength from weekly to six-month period provided more confident strategies compared to other indicators.
    \end{itemize}

    \item An Algorithmic Trading Approach Merging Machine Learning With\\
    Multi-Indicator Strategies for Optimal Performance
    \cite{10795119}
    \begin{itemize}
        \item Investigates the combination of machine learning and multi-indicator strategies for algorithmic trading.
        \item The result showed that their proposed algorithm significantly outperformed traditional indicators.
        \item The research highlights the need of adaptive trading strategies that respond to the dynamics of the market.
    \end{itemize}

    \item Technical indicator empowered intelligent strategies to predict stock trading signals
    \cite{saud2024technical}
    \begin{itemize}
        \item There exists a big gap between how indicators are applied by technical analysts and machine learning experts.
        \item The study proposes trading strategies and four key insights. (1) optimal lookback-period. (2) GRU networks are superior to LSTM. (3) Trading strategies based on MACD, DMI, and KST indicators outperform traditional trading methods. (4) From the 3 proposed intelligent strategies, the MACD was the safest and most effective. 
    \end{itemize}
    
\end{enumerate}

\section{Overfitting}
    \begin{enumerate}
    
        \item Overfitting and undercomputing in machine learning
        \cite {dietterich1995overfitting}

        \begin{itemize}
        \item Introduces the concepts of overfitting and underfitting in machine learning.
        \item Explains how excessive model complexity may reduce generalization performance.
        \item Highlights the importance of computational capacity and proper validation.
        \end{itemize}
        
        \item The problem of overfitting
        \cite {hawkins2004problem}

        \begin{itemize}
        \item Provides a detailed discussion of overfitting in statistical modeling.
        \item Shows how models may fit random noise instead of true structure.
        \item Emphasizes the need for independent test sets and cautious model evaluation.
        \end{itemize}
        
        \item Methods to avoid over-fitting and under-fitting in supervised machine learning (comparative study).
        \cite {jabbar2015methods}

        \begin{itemize}
        \item Compares practical techniques to mitigate overfitting in supervised learning.
        \item Discusses regularization, resampling methods, and model selection approaches.
        \end{itemize}
        
        \item Research on overfitting problem and correction in machine learning.
        \cite {bu2020research}

        \begin{itemize}
        \item Reviews common causes of overfitting in predictive modelling.
        \item Proposes correction strategies to improve model stability and generalization.
        \item Stresses the importance of separating training and testing data properly.
        \end{itemize}
        
        \item Overfitting or poor learning: A critique of current financial applications of GP.
        \cite {shu2003overfitting}

        \begin{itemize}
        \item Critiques the application of genetic programming in financial markets.
        \item Argues that many reported profitable strategies may reflect overfitting.
        \item Highlights the challenge of distinguishing genuine predictive learning from data-specific optimisation in trading.
        \end{itemize}
        
        \item Chaos, overfitting and equilibrium: To what extent can machine learning beat the financial market?.
        \cite {peng2024chaos}

        \begin{itemize}
        \item Investigates whether machine learning models can consistently outperform financial markets.
        \item Examines the relationship between market dynamics and model overfitting.
        \item Suggests that instability in financial systems increases the risk of over-optimised models.
        \item Emphasizes the difficulty of maintaining strong out-of-sample performance.
        \end{itemize}
        
        \item Tackling overfitting in evolutionary-driven financial model induction.
        \cite {tuite2011tackling}

        \begin{itemize}
        \item Focuses on overfitting in evolutionary financial model development.
        \item Proposes techniques to reduce excessive model adaptation to historical data.
        \item Demonstrates the importance of robust evaluation in trading systems.
        \end{itemize}
        
    \end{enumerate}

\section{Validation and Robust Evaluation of Trading Strategies}
    \begin{enumerate}
    
        \item A reality check for data snooping.
        \cite {white2000reality}

        \begin{itemize}
        \item Introduces statistical methods to detect data-snooping bias in empirical research.
        \item Proposes a “reality check” framework for evaluating multiple trading strategies.
        \item Demonstrates how performance may be overstated without correction.
        \end{itemize}
        
        \item Data‐snooping, technical trading rule performance, and the bootstrap.
        \cite{sullivan1999data}

        \begin{itemize}
        \item Examines technical trading rules under data-snooping adjustments.
        \item Uses bootstrap methods to reassess reported profitability.
        \item Shows that many apparent profits disappear after bias correction.
        \end{itemize}
        
        \item The probability of backtest overfitting.
        \cite{bailey2017probability}

        \begin{itemize}
        \item Quantifies the probability of backtest overfitting in trading strategies.
        \item Shows how strategy selection inflates in-sample performance.
        \item Emphasizes the importance of robust out-of-sample testing.
        \end{itemize}
        
        \item Bagging exponential smoothing methods using STL decomposition and Box–Cox transformation.
        \cite{bergmeir2016bagging}

        \begin{itemize}
        \item Discusses validation techniques for time-series forecasting models.
        \item Highlights limitations of traditional cross-validation for dependent data.
        \item Supports the need for proper sequential evaluation procedures.
        \end{itemize}
        
    \end{enumerate}

\section{AI trading}
    \begin{enumerate}
    
        \item Artificial intelligence techniques in financial trading: A systematic literature review
        \cite {dakalbab2024artificial}

        \begin{itemize}
        \item Provides a systematic literature review of AI techniques in financial trading.
        \item Categorizes machine learning approaches used in trading systems.
        \item Identifies challenges related to overfitting and model robustness.
        \end{itemize}
        
        \item Using genetic algorithms to find technical trading rules
        \cite {allen1999using}

        \begin{itemize}
        \item Uses genetic algorithms to discover technical trading rules.
        \item Evaluates whether evolutionary methods can outperform standard
        benchmarks.
        \item Demonstrates early integration of AI techniques in trading.
        \end{itemize}
        
        \item On the profitability of technical trading rules based on artificial neural networks:: Evidence from the Madrid stock market.
        \cite {fernandez2000profitability}

        \begin{itemize}
        \item Applies artificial neural networks to technical trading strategies.
        \item Evaluates profitability in the Madrid stock market.
        \item Shows mixed empirical results for AI-based trading models.
        \end{itemize}
        
    \end{enumerate}

\section{Stop-loss and take-profit}
    \begin{enumerate}
        \item When do stop-loss rules stop losses?
        \cite{kaminski2014stop}
        \begin{itemize}
            \item Framework for measuring the performance of stop-loss rules.
            \item Stop-loss seem effective only if price follows momentum.
        \end{itemize}

        \item Take Profit and Stop Loss Trading Strategies Comparison in Combination with an MACD Trading System
        \cite{vezeris2018take}
        \begin{itemize}
            \item Explores stop-loss and take-profit strategies using MACD.
            \item hard to improve just a simple MACD stop-loss and take-profit strategy.
        \end{itemize}

        \item Effectiveness of Stop-Loss Trading Strategy VS Buy-And-Hold Strategy
        \cite{teimoori2023effectiveness}
        \begin{itemize}
            \item Their testing showed that the trailing stop-loss strategy had the highest return of their strategies.
            \item Trailing stop-loss seems more efficient than buy and hold strategy according to their research.
        \end{itemize}

        \item Optimizing the Pairs-Trading Strategy Using Deep Reinforcement Learning with Trading and Stop-Loss Boundaries
        \cite{kim2019optimizing}
        \begin{itemize}
            \item Investigates the use of machine-learning to optimize stop-loss for trading pairs.
            \item According to their testing, this seems positive for optimizing profits and strategies.
        \end{itemize}
        
    \end{enumerate}

\section{Probability}
    \begin{enumerate}
    
        \item Decision-making for stock trading based on trading probability by considering whole market movement
        \cite{huang2004decision}
        \begin{itemize}
            \item Taking the stock market movement as a whole into consideration when trading individual stocks.
            \item Their testing demonstrates that this method leads with high certainty to better profits than isolated stock strategies.
        \end{itemize}

        \item Intelligent stock trading system by turning point confirming and probabilistic reasoning
        \cite{bao2008intelligent}
        \begin{itemize}
            \item Combines technical analysis with probabilistic reasoning to reduce noise and uncertainty.
        \end{itemize}
        
    \end{enumerate}

\section{Trends and momentum}
    \begin{enumerate}
        \item Stock trading rule discovery with an evolutionary trend following model
        \cite{hu2015stock}
            \begin{itemize}
                \item Introduces a trend following model which can be used to further define strategies based on trends in the market.
                \item based on their research, their eTrend model outperforms other trading models.
                \item outperforms buy-and-hold strategies based on their testing.
            \end{itemize}

        \item Leveraging the momentum effect in machine learning-based cryptocurrency trading
        \cite{bellocca2022leveraging}
        \begin{itemize}
            \item Using machine learning to detect the momentum effect in cryptocurrency to predict short-term volatility and to reduce false trading signals.
        \end{itemize}

        \item Stock Market Trend Analysis Using Hidden Markov Models
        \cite{kavitha2013stock}
        \begin{itemize}
            \item Analysing trends in the stock market using hidden Markov models.
            \item Seems that it might work, but no real testing was done.
        \end{itemize}

        \item Stock Market Trend Analysis and Machine Learning-based\\
        Predictive Evaluation
        \cite{hurriyati2023stock}
        \begin{itemize}
            \item Analyses the stock market with a simple feed-forward neural network model, and contrast Elman, fuzzy logic, and radial basis function networks.
            \item Each technique bring each own advantages, and the results seem positive in their testing, but future research is needed.
        \end{itemize}

        \item Stock Market Trend Prediction Using High-Order Information of Time Series
        \cite{wen2019stock}
        \begin{itemize}
            \item Seeks to answer how to use recent transaction data in order to predict the ups and downs of the stock market.
            \item They aim to filter out noise by leveraging motifs and neural networks.
            \item In their testing, this led to a 4\%-7\% accuracy improvement compared to traditional signal processing methods.
        \end{itemize}
        
    \end{enumerate}

\section{Citing sources other than scientific papers}

You can also cite sources that are not scientific papers. For books, theses, etc the process of searching for bibliographical data is the same as in the case of papers. For websites, a portable choice is the \verb!@misc! bibliographic category together with a \verb!\url! in the \verb!howpublished! field. Please note that the best in this case is to cite a permanently archived version of the website. Such versions can be created using the website \url{web.archive.org}. For example, a description of the kernel of Windows NT can be found in the reference \cite{WinNT}.

\begin{lstlisting}
@misc{WinNT,
  author       = {{808 MultiMedia LLC}},
  title        = {{MS Windows NT} Kernel Description},
  year         = {1999},
  howpublished = {\url{http://web.archive.org/web/20080207010024/http://www.808multimedia.com/winnt/kernel.htm}},
  note         = {Accessed 2010-09-30}
}
\end{lstlisting}

\begin{itemize}
\item \url{https://www.verylongwebsite.com/this/is/a/very/long/url/that/needs/to/break}
\end{itemize}
